\section{Dojikko: неуклюжесть как безоружность и источник симпатии}

Dojikko — неуклюжесть как амплуа, а не случайная черта: персонаж систематически спотыкается, забывает, роняет, путает, и именно системность делает это амплуа, а не просто "неловкой сценой". Юи Хиракасава из той же "K-On!" — прямая противоположность Нодоке из 4.5: где meganekko держит дистанцию наблюдателя, dojikko вообще не держит дистанции, она погружена в момент настолько, что не замечает базовых препятствий вроде забытой дома гитары в день выступления.

Функция dojikko в сцене — обезоруживание через уязвимость: неуклюжесть снижает угрозу, которую мог бы нести персонаж с другими, более острыми чертами, и делает симпатию к нему почти автоматической — трудно всерьёз злиться на того, кто уже сам создаёт себе проблемы без чужой помощи.

Здесь стоит указать на редкий случай, когда сам баг и есть архетип, а не контраст к нему: в отличие от 4.1--4.5, где баг стоит рядом с силой и создаёт объём через контраст, у dojikko неуклюжесть не изъян компетентного персонажа, а прямое содержание роли. Риск в обратную сторону: если dojikko исключительно мила и ни в чём другом не сильна, персонаж рассыпается в чистый moe-элемент (1.5, 2.1) без функции.

Для Рин dojikko-амплуа впрямую не подходит: её резкость и техническая точность (0.4) плохо сочетаются с систематической неловкостью. Но микро-версия полезна ровно как противовес: если Рин никогда не путает, не забывает, не спотыкается даже метафорически в разговоре, она рискует стать той самой безупречной функцией, о которой предупреждают 4.2 и 4.3, — редкая, честная, непостановочная растерянность держит её персонажем, а не движком правильных ответов.
